\subsubsection*{Resultados del SVM}

En este apartado vamos a mostrar los resultados y comentaremos un poco  al respecto, aunque ahondaremos más en la parte de IA explicable (XAI).

\subsubsection*{Modelo baseline con kernel lineal}

============================================================
 EVALUACIÓN: SVM - Lineal
============================================================

MÉTRICAS PRINCIPALES:
   • Accuracy:  1.0000  (Porcentaje total de aciertos)
   • Precision: 1.0000  (De los que dije que eran reales, ¿cuántos lo eran?)
   • Recall:    1.0000  (De los reales, ¿cuántos detecté?)
   • F1-Score:  1.0000  (Balance precision-recall)
   • AUC-ROC:   1.0000  (Capacidad discriminativa)
   • MCC:       1.0000  (Coeficiente de correlación de Matthews)

CLASSIFICATION REPORT:
              precision    recall  f1-score   support

    Real (1)     1.0000    1.0000    1.0000        10
      IA (0)     1.0000    1.0000    1.0000        10

    accuracy                         1.0000        20
   macro avg     1.0000    1.0000    1.0000        20
weighted avg     1.0000    1.0000    1.0000        20


ANÁLISIS DE ERRORES:
   • Verdaderos Negativos (Reales bien clasificados): 10
   • Verdaderos Positivos (IA bien clasificados): 10
   • Falsos Positivos (Reales clasificados como IA): 0
   • Falsos Negativos (IA clasificados como Reales): 0
   • Error tipo I (Falsos Positivos): 0.00%
   • Error tipo II (Falsos Negativos): 0.00%

============================================================ \\

Podemos ver que, pese a ser un modelo baseline, y tener un kernel lineal (que suelen ser peores que los no lienales) ya funciona perfectamente. Esto puede deberse a que los audios, cuando se escuchan, se detectan algunos donde no se entiende lo que está hablando el locutor y otros donde pese a entenderse mejor, por lo que hasta un modelo simple como este ya clasifica bien. Vamos a ver si el kernel no lineal consigue los mismos resultados.

\subsubsection*{Modelo baseline con kernel no lineal}

============================================================
 EVALUACIÓN: SVM - RBF
============================================================

MÉTRICAS PRINCIPALES:
   • Accuracy:  1.0000  (Porcentaje total de aciertos)
   • Precision: 1.0000  (De los que dije que eran reales, ¿cuántos lo eran?)
   • Recall:    1.0000  (De los reales, ¿cuántos detecté?)
   • F1-Score:  1.0000  (Balance precision-recall)
   • AUC-ROC:   1.0000  (Capacidad discriminativa)
   • MCC:       1.0000  (Coeficiente de correlación de Matthews)

CLASSIFICATION REPORT:
              precision    recall  f1-score   support

    Real (1)     1.0000    1.0000    1.0000        10
      IA (0)     1.0000    1.0000    1.0000        10

    accuracy                         1.0000        20
   macro avg     1.0000    1.0000    1.0000        20
weighted avg     1.0000    1.0000    1.0000        20


ANÁLISIS DE ERRORES:
   • Verdaderos Negativos (Reales bien clasificados): 10
   • Verdaderos Positivos (IA bien clasificados): 10
   • Falsos Positivos (Reales clasificados como IA): 0
   • Falsos Negativos (IA clasificados como Reales): 0
   • Error tipo I (Falsos Positivos): 0.00%
   • Error tipo II (Falsos Negativos): 0.00%

============================================================ \\

Vemos que sí, ambos modelos baseline obtienen unos resultados perfectos. Vamos a ver la validación cruzada sobre el modelo de kernel no lineal.

\subsubsection*{Validación Cruzada Estratificada}

Resultados Validación Cruzada SVM RBF (5 folds):
   • Accuracy: 1.0000 (+/- 0.0000)
   • F1-Score: 1.0000 (+/- 0.0000)
   • AUC-ROC:  1.0000 (+/- 0.0000)
   
Estos resultados nos dicen que el modelo funciona bien en los 5 folds, aunque no podemos descartar que hubiera algún fold más ``difícil'' en el cual no tuviera 100\% de precisión.

\subsubsection*{Grid Search}

En primer lugar veamos cuáles son los mejores parámetros del grid que hemos definido.

C: 1
gamma: 0.01
kernel: rbf

Vemos que difieren de los que nosotros definimos previamente en el modelo baseline, y tiene sentido pues la máquina, tras entrenar con todas las posibilidades, ha evaluado este conjunto de parámetros como el mejor, maximizando f1 como vimos. Veamos entonces qué resultados obtenemos con estos parámetros, aunque viendo el resultado del modelo anterior, nos hacemos una idea.

============================================================
 EVALUACIÓN: SVM - Optimizado (RBF)
============================================================

MÉTRICAS PRINCIPALES:
   • Accuracy:  1.0000  (Porcentaje total de aciertos)
   • Precision: 1.0000  (De los que dije que eran reales, ¿cuántos lo eran?)
   • Recall:    1.0000  (De los reales, ¿cuántos detecté?)
   • F1-Score:  1.0000  (Balance precision-recall)
   • AUC-ROC:   1.0000  (Capacidad discriminativa)
   • MCC:       1.0000  (Coeficiente de correlación de Matthews)

CLASSIFICATION REPORT:
              precision    recall  f1-score   support

    Real (1)     1.0000    1.0000    1.0000        10
      IA (0)     1.0000    1.0000    1.0000        10

    accuracy                         1.0000        20
   macro avg     1.0000    1.0000    1.0000        20
weighted avg     1.0000    1.0000    1.0000        20


ANÁLISIS DE ERRORES:
   • Verdaderos Negativos (Reales bien clasificados): 10
   • Verdaderos Positivos (IA bien clasificados): 10
   • Falsos Positivos (Reales clasificados como IA): 0
   • Falsos Negativos (IA clasificados como Reales): 0
   • Error tipo I (Falsos Positivos): 0.00%
   • Error tipo II (Falsos Negativos): 0.00%

============================================================ \\

Como observamos y como la intuición nos decía, los resultados son iguales que los de antes, todo perfecto. Posiblemente la razón sea la misma: una mezcla entre un buen split y pocos datos y/o datos poco representativos.